% 参数积分中质量逃逸：f(t,x)=t e^{-tx}，t=2,1,0.5,0.25，每条曲线下面积为 1；虚线为包络 sup_t f = 1/(e x)（在 t=1/x 取到），在 [1,8] 上画出，它在无穷远处不可积。
\begin{tikzpicture}[>=Stealth,
  ax/.style={draw=FigureStroke,line width=.6pt,->},
  note/.style={font=\FigureNoteFont\fontsize{7.5}{9}\selectfont,text=FigureMuted}]
\begin{scope}[xscale=1.1,yscale=2.0]
\draw[FigureGrid,line width=.4pt] (0,.5)--(8,.5) (0,1)--(8,1) (0,1.5)--(8,1.5);
\draw[ax] (-.1,0)--(8.3,0) node[right] {$x$};
\draw[ax] (0,-.05)--(0,2.15) node[above] {$f(t,x)$};
\foreach \y in {0.5,1,1.5,2}{\node[left,font=\fontsize{7.5}{9}\selectfont] at (0,\y) {$\y$};}
\foreach \x in {2,4,6,8}{\node[below,font=\fontsize{7.5}{9}\selectfont] at (\x,0) {$\x$};}
\draw[FigureBlue,line width=1.1pt,domain=0:8,samples=80] plot (\x,{2*exp(-2*\x)});
\draw[FigureGreen,line width=1.1pt,domain=0:8,samples=80] plot (\x,{exp(-\x)});
\draw[FigureYellow,line width=1.1pt,domain=0:8,samples=80] plot (\x,{.5*exp(-.5*\x)});
\draw[FigureRed,line width=1.1pt,domain=0:8,samples=80] plot (\x,{.25*exp(-.25*\x)});
\draw[FigureStroke,dashed,line width=.7pt,domain=1:8,samples=60] plot (\x,{1/(2.71828*\x)});
\node[text=FigureBlue,anchor=west] at (.25,1.75) {$t=2$};
\node[text=FigureGreen,anchor=west] at (.7,.7) {$t=1$};
\node[text=FigureYellow,anchor=west] at (2.0,.33) {$t=0.5$};
\node[text=FigureRed,anchor=west] at (6.0,.15) {$t=0.25$};
\node[anchor=west,font=\fontsize{7.5}{9}\selectfont] at (4.4,.22) {包络$1/(ex)$};
\end{scope}
\end{tikzpicture}
